\let\negmedspace\undefined
\let\negthickspace\undefined
\documentclass[journal]{IEEEtran}
\usepackage[a5paper, margin=10mm, onecolumn]{geometry}
\usepackage{tfrupee} 

\setlength{\headheight}{1cm} 
\setlength{\headsep}{0mm}     
\usepackage{gvv-book}
\usepackage{gvv}
\usepackage{cite}
\usepackage{amsmath,amssymb,amsfonts,amsthm}
\usepackage{algorithmic}
\usepackage{graphicx}
\usepackage{textcomp}
\usepackage{xcolor}
\usepackage{txfonts}
\usepackage{listings}
\usepackage{enumitem}
\usepackage{mathtools}
\usepackage{gensymb}
\usepackage{comment}
\usepackage[breaklinks=true]{hyperref}
\usepackage{tikz}
\usetikzlibrary{arrows.meta,angles,quotes}

\begin{document}

\bibliographystyle{IEEEtran}
\vspace{3cm}

\title{5.3.2}
\author{AI25BTECH11008 - Chiruvella Harshith Sharan}
{\let\newpage\relax\maketitle}

\setlength{\intextsep}{10pt} 
\numberwithin{equation}{enumi}

\textbf{Question:} \\

\vspace{0.3cm}

The pair of linear equations 
\[
\frac{3x}{2} + \frac{5y}{3} = 7 
\quad \text{and} \quad 
9x + 10y = 14
\]
is 

\begin{enumerate}[label=(\alph*)]
\item consistent 
\item inconsistent 
\item consistent with one solution 
\item consistent with many solutions
\end{enumerate}

\vspace{0.3cm}

\textbf{Solution:}

\subsection*{Step 1: Clear fractions and write augmented matrix}

Multiply the first equation by $6$ to clear fractions, obtaining $9x+10y=42$. We now have
\begin{align}
9x + 10y &= 42 \tag{1}\\
9x + 10y &= 14 \tag{2}
\end{align}
The corresponding augmented matrix is
\begin{align}
\begin{pmatrix}
9 & 10 & \vline & 42\\[2pt]
9 & 10 & \vline & 14
\end{pmatrix}. \tag{3}
\end{align}

\subsection*{Step 2: Row operation to check consistency}

Perform the row operation $R_2 \leftarrow R_2 - R_1$ to eliminate the coefficients:
\begin{align}
\begin{pmatrix}
9 & 10 & \vline & 42\\[2pt]
0 & 0  & \vline & -28
\end{pmatrix}. \tag{4}
\end{align}

The second row corresponds to $0x+0y=-28$, which is impossible. Hence the system has no solution.

\subsection*{Final Answer:}

\[
\boxed{\text{The system is inconsistent.}}
\]

\begin{figure}[h!]
    \centering
    \includegraphics[width=0.6\linewidth]{figs/fig1.png}
    \caption{Graph of the lines $9x+10y=42$ and $9x+10y=14$.}
    \label{fig:lines}
\end{figure}


\end{document}
